\section{Final Remarks}
\label{sec:final}

We introduced tableau calculi for data-aware logics, with particular attention to variants of the query language XPath extended with data comparisons. Precisely, we presented an internalized tableau calculus for the logic called $\Hxpath$. This logic extends the standard XPath formalism as studied in the context of Modal Logic by incorporating data comparisons, along with nominals and satisfiability modalities from Hybrid Logic.
We proved that our tableau calculus is sound, complete, and terminating, making it a practical tool for checking the satisfiability of $\Hxpath$ node and path expressions. Additionally, we showed that rule applications can be systematically restricted to ensure that tableau construction runs in polynomial space, without sacrificing completeness. This yields a \pspace-completeness result for the satisfiability problem in $\Hxpath$.

Our work builds upon~\cite{ArecesFS17} and offers several enhancements. Most notably, it provides a detailed account of the termination argument, which was only briefly outlined in the earlier article. The presence of nominals and satisfiability modalities in the language implies that the tableau calculi need to perform equality reasoning, which complicates termination. In particular, the introduction of new nominals needs to be carefully handled.  In addition, we demonstrate that the \textsc{PSpace} upper bound also applies in the case of trees.
The current proposal is, arguably, simpler and more elegant than the presentation of the calculus given in~\cite{ArecesFS17}.  In particular, the \pspace upper bound for satisfiability is presented in full, at a level of detail suitable for implementation. The tableau calculi introduced in this article are \emph{internalized}, i.e., rules only contain expressions of $\Hxpath$. It is known that internalized proof systems for Hybrid Logics enjoy some nice properties~\cite{Blackburn00,Blackburn2002BeyondPA}. We exploited these properties here, demonstrating how to extend the tableau calculus with pure axioms and node creating rules to characterize interesting model classes. Moreover, we showed that strong completeness theorems holds for these particular kind of tableau rules, ensuring that the resulting calculi are automatically complete for the corresponding model classes.  We illustrated how these rules can be used to capture new operators and model classes relevant for applications based on semi-structured data. 

%The contributions of this paper are twofold. On the one hand, we contributed to the understanding of XPath from a Modal Logic perspective. On the other hand, we established a first step towards studying Modal Logics extended with data features in general. 

The results presented in this article pave the way to a modal investigation of other data-aware logics that are receiving considerable attention recently, such as GQL~\cite{FrancisGGLMMMPR23,FrancisGGLMMMPR23b} and SHACL~\cite{Ortiz23,Ahmetaj0S23}. It also makes it possible to consider Modal Logics built over more complex data structures, such as Concrete Domains~\cite{DemriQ21}. It would be interesting to explore futher connections with Register Automata~\cite{KaminskiF94} and similar structures. Finally, implementing a tableaux-based procedure for data-aware languages over the hybrid prover \textsf{HTab}~\cite{Hoffmann2007} is also part of our agenda.
